% /Users/pikachu/books/payments-book/src/parts/part5/ch27-multicurrency.tex
\chapter{Multi-currency и cross-border}
\label{ch:multicurrency}
\margintoc


\begin{kaobox}[title=Что вы узнаете из этой главы]
  Когда покупатель в Москве платит картой российского банка
  на сайте европейского мерчанта, в транзакции участвуют три
  валюты: валюта транзакции, валюта биллинга и валюта settlement.
  Кто конвертирует, на каком этапе, кто теряет на курсе,
  что такое DCC и почему от него лучше отказаться,
  как устроены трансграничные расчёты через корреспондентские
  цепочки и что меняет ISO 20022.
\end{kaobox}

Возьмём одну транзакцию: держатель рублёвой карты в Москве платит
100 EUR европейскому мерчанту. Почти все дальнейшие проблемы —
курс списания, DCC, FX risk, settlement delays — становятся
понятными, если с самого начала не смешивать три валюты и три
разных места, где они фиксируются.

\section{Три валюты одной транзакции}
\label{sec:mc-currencies}

Внутри одной страны валютный вопрос не возникает: покупатель
платит рублями, мерчант получает рубли, и единственная конверсия —
из копеек в рубли при отображении суммы на экране. Но как только
транзакция пересекает границу, в ней появляются до трёх разных
валют, и понимание того, какая из них где фиксируется и кем
конвертируется, критично для корректного учёта и управления
валютным риском.

\begin{kaobox}[title={Определение: Transaction, billing и settlement currency}, colback=blue!4, colbacktitle=blue!15]
  \textbf{Transaction currency} — валюта, в которой мерчант
  выставляет цену и в которой номинирована транзакция;
  фиксируется в \de{49} (Currency Code, Transaction)
  сообщения \iso{8583}.
  \textbf{Billing currency} — валюта, в которой эмитент
  списывает средства со счёта держателя карты; определяется
  валютой карточного счёта.
  \textbf{Settlement currency} — валюта, в которой карточная
  сеть рассчитывается с эквайером (и через него — с мерчантом);
  фиксируется в \de{50} (Currency Code, Settlement).
\end{kaobox}

Рассмотрим пример. Российский держатель рублёвой карты
\scheme{Visa} покупает товар за 100 EUR на сайте немецкого
мерчанта. Transaction currency — EUR (\de{49} = 978),
потому что мерчант выставил цену в евро. Settlement
currency — тоже EUR, потому что эквайер мерчанта работает
в еврозоне и получает settlement в евро. Billing currency —
RUB, потому что карточный счёт держателя номинирован в рублях.

Конверсия EUR → RUB происходит на стороне эмитента (или карточной
сети от имени эмитента) по курсу, установленному \scheme{Visa}
на дату клиринга — не на дату авторизации. Это означает,
что держатель карты в момент покупки не знает точную сумму
списания в рублях: он видит авторизационный hold,
рассчитанный по курсу на дату авторизации, но фактическое
списание произойдёт позже, по другому курсу, и сумма может
отличаться — как в большую, так и в меньшую сторону.

\begin{kaobox}[title=Важно, colback=red!5, colbacktitle=red!20]
  Разрыв между курсом авторизации и курсом клиринга —
  источник одного из самых частых вопросов в поддержку:
  \enquote{мне списали больше, чем было на экране}. Технически
  это не ошибка, а следствие двухфазной модели: авторизация
  фиксирует сумму в transaction currency, а конверсия
  в billing currency происходит на этапе клиринга, когда
  курс мог измениться. Для волатильных валютных пар
  (рубль/евро, турецкая лира/доллар) разница может
  достигать нескольких процентов.
\end{kaobox}

Итак, как только transaction, billing и settlement currency разведены
по местам, становится видно, где именно возникает неопределённость
для держателя и валютный риск для участников цепочки. DCC —
это первая попытка эту неопределённость монетизировать,
а не устранить бесплатно.

\section{DCC: Dynamic Currency Conversion}
\label{sec:mc-dcc}

DCC — механизм, при котором конверсия валюты происходит не на стороне эмитента (как в стандартном flow), а на стороне
эквайера: терминал или платёжная страница предлагает держателю карты оплатить покупку не в валюте мерчанта, а сразу в
валюте его карточного счёта, показывая курс конверсии и итоговую сумму.

На первый взгляд DCC удобен: покупатель видит сумму в знакомой
валюте и точно знает, сколько будет списано, — устраняя
неопределённость, описанную в предыдущем разделе. Но удобство
это дорого обходится: DCC-провайдер (обычно эквайер
или специализированная компания) закладывает markup — наценку
к оптовому курсу (wholesale rate), типично составляющую
3--7\%, что значительно превышает markup карточной
сети \todo[inline]{Проверить: стандартный FX markup Visa/MC для cross-border
  конверсии — обычно около 1\%; сравнить с типичным DCC
  markup 3--7\%}.

Правила \scheme{Visa} и \scheme{Mastercard} требуют,
чтобы DCC предлагался только как opt-in: держатель карты должен
явно согласиться на конверсию, видя полный курс и markup.
Принудительное применение DCC (без явного согласия) —
нарушение правил сети, за которое эквайер может быть
оштрафован.

\begin{kaobox}[title=На практике]
  Если вы разрабатываете платёжный flow для мерчанта
  с международным трафиком — не включайте DCC по умолчанию.
  Для большинства держателей карт конверсия по курсу карточной
  сети (без DCC) выгоднее, чем DCC с markup эквайера.
  DCC имеет смысл предлагать только в странах, где курс
  карточной сети аномально невыгоден или где покупатель
  особенно ценит предсказуемость суммы (например,
  корпоративные карты с жёсткими бюджетами).
\end{kaobox}

С технической точки зрения DCC меняет routing конверсии:
вместо стандартной схемы \enquote{эквайер получает settlement
  в transaction currency, эмитент конвертирует в billing currency}
транзакция приходит к эмитенту уже в billing currency —
конверсия произошла на стороне эквайера до авторизации.
Это означает, что эмитент не применяет свой FX rate,
и весь FX margin остаётся у DCC-провайдера. Именно
это экономическое перераспределение делает DCC привлекательным
для эквайеров и невыгодным для держателей карт.

\section{FX risk и модели расчёта}
\label{sec:mc-fx-risk}

Валютный риск (FX risk) в карточных транзакциях возникает
из-за временного лага между авторизацией и settlement:
курс конверсии фиксируется на дату клиринга, а не на дату
авторизации, и за время между ними курс может измениться.

Для эмитента FX risk проявляется в том, что hold
на карте (рассчитанный по курсу авторизации) может не совпадать
с фактическим списанием (по курсу клиринга). Эмитенты
обычно закладывают буфер: hold рассчитывается по курсу
с небольшим запасом (1--3\%), чтобы покрыть возможное
изменение курса.

Для мерчанта FX risk определяется моделью settlement.
Существуют два основных варианта:

\textbf{Single-currency settlement} — мерчант получает
все средства в одной валюте (обычно — валюте своей страны),
независимо от валюты транзакции. Конверсию выполняет
эквайер или карточная сеть по своему курсу. Мерчант не несёт
прямого FX risk (он всегда получает рубли/евро/доллары),
но теряет на FX markup, заложенном в курс конверсии.

\textbf{Multi-currency settlement} — мерчант получает средства
в валюте транзакции: если покупатель заплатил в EUR, мерчант
получает EUR; если в USD — USD. Мерчант избегает двойной
конверсии (transaction currency → settlement currency →
валюта мерчанта), но берёт на себя FX risk: теперь у него
на счетах несколько валют, и управление их балансом —
его задача.

\begin{kaobox}[title=На практике]
  Для мерчанта с мультивалютным трафиком выбор между single-
  и multi-currency settlement — trade-off между простотой
  и стоимостью. Если основной трафик — в одной валюте
  (например, 90\% в рублях и 10\% в евро), single-currency
  settlement проще операционно: один счёт, один баланс,
  один отчёт. Если трафик распределён между несколькими
  валютами примерно поровну, multi-currency settlement
  экономит на конверсии и позволяет платить поставщикам
  в валюте без обратной конверсии.
\end{kaobox}

\section{Корреспондентский banking и трансграничные расчёты}
\label{sec:mc-correspondent}

В главе \ref{ch:banks} мы разобрали, как работают
корреспондентские счета для обычных межбанковских переводов.
В cross-border карточных транзакциях тот же механизм
используется на этапе settlement, но с существенным отличием:
между банками стоит карточная сеть, которая выступает
расчётным посредником.

При внутристрановых транзакциях settlement проходит через
национальную инфраструктуру: в России — через НСПК
и платёжную систему Банка России. При cross-border —
через расчётные банки карточных сетей: \scheme{Visa}
использует JPMorgan Chase как основной settlement bank,
\scheme{Mastercard} — сеть расчётных банков
в разных регионах \todo[inline]{Источник: Visa и Mastercard settlement
  banks — проверить актуальный перечень расчётных банков}.

Цепочка cross-border settlement выглядит следующим образом.
Карточная сеть рассчитывает нетто-позиции всех участников:
сколько каждый эквайер должен получить и сколько каждый
эмитент должен заплатить, с учётом interchange и комиссий.
Затем сеть инициирует проводки через расчётный банк:
банк-эмитент перечисляет средства со своего корреспондентского
счёта в расчётном банке на корреспондентский счёт
банка-эквайера — или через цепочку банков-корреспондентов,
если прямого корреспондентского отношения нет.

\todo[inline]{Рисунок: Цепочка трансграничного карточного settlement:
  карточная сеть (в центре) рассчитывает нетто-позиции;
  стрелки от эмитентов (слева) через settlement bank
  к эквайерам (справа); пунктиром — корреспондентская
  цепочка для банков без прямого доступа к settlement bank.
  Горизонтальный flow, три колонки: эмитенты, сеть + settlement
  bank, эквайеры.}

\subsection*{SWIFT и миграция на ISO 20022}

Инструкции о переводе средств между банками-корреспондентами
исторически передавались через SWIFT в формате MT:
MT103 для клиентских переводов и MT202 для межбанковских.

С 2020-х годов отрасль мигрирует на формат MX, основанный
на \iso{20022}. Это международный стандарт финансовых сообщений
с расширяемой XML-схемой, поддержкой структурированных данных
и единой моделью данных для разных типов операций
(см. главу \ref{ch:banks}).

Для платёжного процессинга миграция на \iso{20022} означает
несколько вещей. Во-первых, более богатые данные в settlement-сообщениях:
\iso{20022} позволяет передать structured remittance information
(детали платежа, invoice reference), которые в формате MT
приходилось упаковывать в свободные текстовые поля.
Во-вторых, улучшенная совместимость: \iso{20022} — единый
формат для SWIFT, БЭСП, SEPA, Fedwire и десятков других
платёжных систем, что упрощает маршрутизацию сообщений
через разные инфраструктуры. В-третьих, compliance:
структурированные данные облегчают автоматический screening
(санкционный скрининг, AML-проверки), который для свободных
текстовых полей MT требовал эвристик с высоким уровнем
false positives.

\subsection*{SWIFT GPI}

SWIFT GPI (Global Payments Innovation) — инициатива, запущенная в 2017 году для решения трёх хронических проблем
международных переводов: непрозрачность (отправитель не знает, где находится его платёж), непредсказуемость (время и
стоимость перевода зависят от числа банков-корреспондентов) и отсутствие подтверждения \sidecite{swift-overview}.

GPI вводит end-to-end tracking. Каждому платежу присваивается уникальный UETR (Unique End-to-End Transaction
Reference), который сохраняется по всей цепочке банков-корреспондентов.

Каждый банк в цепочке обязан обновлять статус платежа
в GPI Tracker, и отправитель может в реальном времени видеть,
на каком этапе находится перевод. Дополнительно GPI задаёт SLA:
банк-участник должен зачислить same-currency перевод в тот же
операционный день или обосновать задержку.

\subsection*{СПФС и расчёты в СНГ}

Для внутрироссийских расчётов и расчётов с дружественными
странами всё большую роль играет СПФС — российский аналог
SWIFT, описанный в главе \ref{ch:banks}. После 2022 года,
когда ряд российских банков был отключён от SWIFT, СПФС стала
основным каналом для передачи финансовых сообщений внутри
России и активно подключает банки из стран СНГ, Китая, Индии
и других.

Трансграничные расчёты в СНГ развиваются в нескольких
направлениях. Двусторонние соглашения между центральными
банками (Россия--Казахстан, Россия--Беларусь) позволяют
проводить расчёты в национальных валютах, минуя долларовую
инфраструктуру. Кобейджинговые карты (Мир--UnionPay,
Мир--Белкарт) обеспечивают приём российских карт за рубежом
через партнёрские сети (см. главу \ref{ch:mir}). Интеграция
национальных систем быстрых платежей (СБП — Россия,
Kaspi — Казахстан) находится на стадии обсуждения
и пилотных проектов.

\begin{kaobox}[title=На практике]
  Если вы проектируете cross-border платёжный flow
  для рынков СНГ, учитывайте фрагментированность
  инфраструктуры: для каждой пары стран набор доступных
  rails (карты, A2A, SWIFT/СПФС) и валютных маршрутов
  различается. Единого решения нет, и архитектура должна
  поддерживать конфигурируемый routing по коридору
  (corridor-based routing): для каждой пары
  \enquote{страна отправителя — страна получателя}
  определять оптимальный метод и валюту расчёта.
\end{kaobox}

\section*{Итоги главы}

\begin{itemize}
  \item В cross-border транзакции участвуют три валюты:
        transaction currency (цена мерчанта, \de{49}),
        settlement currency (расчёт с эквайером, \de{50})
        и billing currency (списание с карты); конверсия
        в billing currency происходит по курсу на дату
        клиринга, а не авторизации.
  \item DCC переносит конверсию на сторону эквайера; markup DCC-провайдера (3--7\%) обычно значительно превышает FX markup
        карточной сети ($\approx$1\%); правила сетей требуют явного opt-in от держателя.
  \item Multi-currency settlement экономит на конверсии, но требует управления валютными балансами; single-currency проще
        операционно, но дороже из-за FX markup.
  \item Cross-border settlement проходит через расчётные банки карточных сетей и корреспондентские цепочки; SWIFT мигрирует с
        MT на MX (\iso{20022}); GPI обеспечивает end-to-end tracking и SLA.
  \item В СНГ трансграничные расчёты развиваются через
        СПФС, двусторонние соглашения в национальных валютах
        и кобейджинг; инфраструктура фрагментирована,
        и архитектура должна поддерживать corridor-based
        routing.
\end{itemize}
